%%%%%%%%%%%%%%%%%%%%%%%%%%%%%%%%%%%%%%%%%%%%%%%%%%%%%%%%%%%%%%%%%%
\section{Cross-benchmark analysis}\label{sec:cross_benchmark}
%%%%%%%%%%%%%%%%%%%%%%%%%%%%%%%%%%%%%%%%%%%%%%%%%%%%%%%%%%%%%%%%%%
% Synthesis of Section 5. All values are repeated from Section 5 / the
% evidence inventory / paper-data evidence reports, except the Hessenthaler
% convergence-test evidence (Section iterative_errors), which has no Section-5
% counterpart because that case has no convergence results yet.

\todo{Synthesis based on the evidence available on 2026-10-10. Conclusions will change when: (a) the beamInCrossFlow (Richter) and Hessenthaler verification campaigns are completed; (b) independent references become available (3dTube, mokLidDrivenCavity first); (c) the Hron--Turek geometric-singularity hypothesis is tested.}

This section draws together the findings of Section~\ref{sec:benchmark_suite} by theme.
In this suite the strength of a verification statement is usually limited by the quality of the reference, the number of refinement levels that could be afforded, or the iterative error, or by sub-nominal convergence whose cause is only partly identified.
Verification also exposed defects that agreement with a reference would not have shown: an $O(\Delta t)$ boundary-flux inconsistency (womersleyTube), a compiler-dependent aliasing error and a mesh-dependent interpolation degeneracy (3dTube), and a coupling convergence test that can accept unconverged steps (Hessenthaler); these are collected in Section~\ref{sec:defects}.
The observations are specific to this suite and to one code family, and are not offered as universal.

%%%%%%%%%%%%%%%%%%%%%%%%%%%%%%%%%%%%%%%%%%%%%%%%%%%%%%%%%%%%%%%%%%
\subsection{Nominal and observed convergence}\label{sec:cross_convergence}
%%%%%%%%%%%%%%%%%%%%%%%%%%%%%%%%%%%%%%%%%%%%%%%%%%%%%%%%%%%%%%%%%%

\begin{table}[htbp]
	\centering
	\footnotesize
	\setlength{\tabcolsep}{3pt}
	\caption{Observed orders across the suite (formal order 2 in space and time). ``SD'': successive differences of three or more levels (reference-free); ``vs exact'': error against an exact reference; ``2-level vs ref.'': apparent order from two levels against a non-exact reference; ``scaled'': $\Delta t$ reduced with $h$.}
	\label{tab:orders}
	\begin{tabular}{p{3.0cm}p{4.6cm}p{3.4cm}p{2.8cm}}
		\toprule
		Case & Space & Time & Basis \\
		\midrule
		ringAddedMass & freq.\ converged on coarsest mesh; ratio 1.68--1.77; standard solid 1.66 & 1.90--2.14 & vs exact; Richardson \\
		womersleyTube & 2.01--2.05 (amplitudes, profile, speed); 1.92 (flow phase); 1.40 (wall phase); attenuation at linear-theory level & 1.86--2.14 (1.29--1.84 before end-flux fix) & SD; errors vs exact \\
		collapsibleChannel & 1.7, 2.4, then reference floor & 1.8, 2.0 & vs ref.\ (B); SD (time) \\
		Hron--Turek FSI1 & 1.5--1.7 & -- & 2-level vs ref.\ (B) \\
		Hron--Turek FSI3 & displacements 0.9--1.1; drag amplitude 0.74 (scaled) & $\leq1.7\%$ change (2x, $\Delta t$ and tolerance) & SD \\
		3dTube & $u_{z,\min}$ 0.87, $t_\mathrm{arr}$ 0.72; $u_{r,\max}$ not defined (changes grow) & 0.1--0.2\% change at the fine-level $\Delta t$ & SD \\
		mokLidDrivenCavity & peak $\approx1.2$; trough $<0.5$ (scaled) & changes $0.17$, $0.04\%$ & SD \\
		cavityFlexibleBottom$^\ast$ & 1.96 ($u_y$; one triple, level 4 diverged) & -- & SD \\
		beamInCrossFlow$^\ast$ & in progress & -- & -- \\
		blobInTreacle$^\ast$ & 2.40 (super-nominal) & 2.39--2.54 (super-nominal) & SD \\
		\bottomrule
		\multicolumn{4}{l}{$^\ast$ Provisional cases.}
	\end{tabular}
\end{table}

Table~\ref{tab:orders} collects the observed orders, from which three observations follow.

\emph{Observed orders can depart from the nominal order.}
Among the category-A cases, ringAddedMass exhibits near-second-order temporal convergence, while the amplitude-type womersleyTube quantities exhibit near-second-order spatial convergence.
Independently of reference quality, collapsibleChannel also shows near-second-order behaviour, in space against its category-B reference (orders 1.7 and 2.4 before the reference floor) and in time by self-convergence (1.8 and 2.0).
womersleyTube first showed temporal orders of 1.3--1.8 that fell towards one as $\Delta t$ was refined, on every mesh.
The orders are reference-free observations, but the diagnosis relied on the exact reference and on fluid-only and end-flux diagnostics: the deficit came from an $O(\Delta t)$ mass-flux inconsistency at the artificial tube ends, where the exact pressure and velocity gradient are imposed, not from the coupling (Section~\ref{sec:womersley}).
With a flux-consistent end treatment every QoI converges at second order in time over the steps used (formally, an $O(\Delta t\,h)$ boundary term remains, as for standard fixed-pressure outlets), and the finest solutions lie within $0.12\%$ of the continuum solution ($0.2\%$ for the attenuation, at the linear-theory level).
In ringAddedMass, by contrast, the exact reference allows the remaining $-0.13\%$ frequency error at a practical resolution to be attributed to the BDF2 phase error.

\emph{Two levels do not establish convergence, and super-nominal orders are a warning.}
Third levels changed the conclusions drawn from two in both cases where they were added.
In Hron--Turek FSI3 the two-level apparent orders of 2.6--3.8 arose from a sequence crossing the reference value, and the three-level observed order of the displacements is about one.
In 3dTube the small two-level change in $u_{r,\max}$ reported earlier was obtained with an interface interpolation operator that changed between the levels; after correction its changes grow on the third level.
FSI1 and FSI2 still have two levels, so the most widely used benchmarks in the suite remain the least completely verified.
Orders well above the formal value (the FSI3 filtered lift amplitude, 3.1; the rigid-flag force amplitudes, 2.8--3.3; the blobInTreacle temporal orders, 2.4--2.5) are interpreted as pre-asymptotic; no super-convergence is claimed, and they are not used for extrapolation.

\emph{Coupled limit-cycle QoIs can converge sub-nominally when the components behave well.}
In Hron--Turek FSI3 the solid alone approaches second order, the rigid-flag fluid converges to within the reference uncertainty, and the code defects of 3dTube have no effect; yet the coupled displacements converge at an observed order of about one, and the in-phase fluid force under the replayed coupled motion shows the same behaviour under every valid fluid treatment tested (Section~\ref{sec:fsi23}).
A self-excited amplitude is set by a balance of small differences between large loads, so a sub-nominal order in the in-phase load can appear as a sub-nominal order in the amplitude.
Two statements must therefore be kept apart: ``our result disagrees with the reference'' and ``our result is not yet resolved''.
At 4x the FSI3 displacements differ from Featflow by more than the reference's level-to-level change, but they are still changing at first order; neither the code nor the reference is shown to be wrong by this comparison.

\emph{Combined space--time refinement conceals the balance of errors.}
Where space and time have been separated, the balance differs strongly: in 3dTube the temporal error at the fine-level $\Delta t$ ($0.1$--$0.2\%$) is several times smaller than the level-2-to-3 change ($0.3$--$1.0\%$), in mokLidDrivenCavity it is an order of magnitude below the spatial error, in ringAddedMass it dominates the frequency, and in womersleyTube, once the end-flux inconsistency is removed, it is a tenth or less of the finest-mesh error for most QoIs (before, the two were similar).
Scaled studies cannot show which dominates unless supplemented, as for Hron--Turek FSI3 and 3dTube, by time-step checks on the finer levels.

%%%%%%%%%%%%%%%%%%%%%%%%%%%%%%%%%%%%%%%%%%%%%%%%%%%%%%%%%%%%%%%%%%
\subsection{Dependence on the quantity of interest}\label{sec:qoi_dependence}
%%%%%%%%%%%%%%%%%%%%%%%%%%%%%%%%%%%%%%%%%%%%%%%%%%%%%%%%%%%%%%%%%%

In every case with more than one QoI, the QoIs converge at substantially different rates (Table~\ref{tab:qoi_dependence}), so a verification statement made on one QoI says little about the others.

\begin{table}[htbp]
	\centering
	\footnotesize
	\setlength{\tabcolsep}{3pt}
	\caption{Faster- and slower-converging QoIs within each case (finest available level).}
	\label{tab:qoi_dependence}
	\begin{tabular}{p{2.8cm}p{5.4cm}p{5.4cm}}
		\toprule
		Case & Faster & Slower \\
		\midrule
		womersleyTube & flow-rate, wall amplitude, wave speed: order $\approx2$, error $\leq0.04\%$ & wall phase (order 1.40); attenuation (error $-0.2\%$, at the linear-theory level) \\
		Hron--Turek FSI1 (2x) & drag $0.24\%$ & $u_y$ $4.64\%$ \\
		Hron--Turek FSI3 (4x) & drag mean and frequency: within $0.3\%$ of the reference & displacements: order $\approx1$; force amplitudes $7.6$--$10.4\%$ above the reference, not asymptotic \\
		Hron--Turek FSI2 (2x) & drag mean $0.5\%$; $u_y$ amplitude $1.3\%$ & lift amplitude $7.8\%$, not converging \\
		3dTube & $u_{z,\min}$, $t_\mathrm{arr}$: monotone, orders 0.87, 0.72 & $u_{r,\max}$, $c_p$: changes grow ($u_{r,\max}$ $0.20\%$, $0.60\%$) \\
		mokLidDrivenCavity & peak (order $\approx1.2$) & trough (order $<0.5$) \\
		cavityFlexibleBottom$^\ast$ & $F_y$: within $0.2\%$ on all levels & $u_y$: 13.2\% $\to$ 5.6\% \\
				\bottomrule
	\end{tabular}
\end{table}

The pattern is not simply that displacements converge faster than forces.
Time-mean integrated forces (Hron--Turek drag, cavityFlexibleBottom vertical force) are often the \emph{fastest}-converging quantities, plausibly because they are dominated by a global pressure balance.
The slowest are of three kinds:
(i) oscillation amplitudes of forces that result from near-cancellation of larger contributions, above all the Hron--Turek lift amplitude, which is also the QoI most sensitive to the coupling tolerance;
(ii) phase- and attenuation-type quantities, or small components, such as the Womersley wall phase and attenuation ($\mathrm{Im}(k)$ is 14\% of $\mathrm{Re}(k)$) and the transverse tip displacement of FSI1; and
(iii) loads that set a self-excited response, such as the in-phase fluid force on the Hron--Turek flag, which converges at an observed order of about one under a fixed motion although the force amplitude converges quickly.
A benchmark should therefore report several QoI types (a displacement, a mean force, an oscillatory force amplitude or phase, and a frequency where applicable), and the slowest QoI on which a claim is made determines the resolution required.

%%%%%%%%%%%%%%%%%%%%%%%%%%%%%%%%%%%%%%%%%%%%%%%%%%%%%%%%%%%%%%%%%%
\subsection{Iterative and coupling errors}\label{sec:iterative_errors}
%%%%%%%%%%%%%%%%%%%%%%%%%%%%%%%%%%%%%%%%%%%%%%%%%%%%%%%%%%%%%%%%%%

The iterative error is the least systematically quantified component of the error budget in the present suite.

\emph{Direct tolerance evidence is limited.}
Hron--Turek FSI2 contains a coupling-tolerance sweep, and it shows that a tolerance of $10^{-5}$ inflates the lift amplitude by $13\%$ (1x) and $30\%$ (2x): the effect grows with refinement, so a tolerance chosen on a coarse mesh can contaminate a fine-mesh result.
FSI3 shows the same mechanism in time: at a fixed tolerance of $10^{-5}$ the step-to-step force noise grows as $\Delta t$ falls (lift $7.3\,\mathrm{N/m}$ rms at the 4x step against $1.7\,\mathrm{N/m}$ at the 2x step), inflating extrema-based force amplitudes while leaving the displacements unaffected, and at the finest step IQN-ILS stalls above a tolerance of $10^{-6}$.
Elsewhere, the iterative error is bounded only indirectly.

\emph{A converged-looking QoI can carry an accumulated coupling error.}
In the Hessenthaler verification campaign (Section~\ref{sec:hessenthaler}; OpenFOAM-v2412, IQN-ILS, tolerance $10^{-4}$, $\Delta t = 0.002\,\mathrm{s}$, 40\,000 steps to $t = 80\,\mathrm{s}$), the legacy test $\min(r_1, r_2)$ of Eq.~\eqref{eq:coupling_residuals} accepted steps after about three iterations once the flow had become quasi-steady, because $r_2$, normalised by the total tip deflection of $\approx17\,\mathrm{mm}$, was small (median $3.5\times10^{-5}$) whatever the reduction within the step.
After $t = 1\,\mathrm{s}$ every step was accepted with $r_1 > 10^{-4}$ (median $\approx6.8\times10^{-2}$; $87$--$91\%$ of steps with $r_1 > 10^{-2}$), the fluid interface lagged the solid by $\approx0.03\,\mathrm{mm}$, and the flap tip crept upwards at $+0.003$--$0.005\,\mathrm{mm/s}$ ($16.84$, $16.90$ and $16.98\,\mathrm{mm}$ at $t = 30$, $60$ and $80\,\mathrm{s}$), by more than the spatial differences under study.
Neither the fluid nor the solid alone drifted.
Restarted from the $80\,\mathrm{s}$ state with the strict test $\max(r_1, r_2)$ (10--11 iterations per step), the tip settled within $\approx0.3\,\mathrm{s}$ to a plateau of $17.0195\,\mathrm{mm}$, flat to $\pm0.0002\,\mathrm{mm}$ over $7.5\,\mathrm{s}$.
Strict restarts from different default-test states did not return to one value ($17.02$ and $16.90\,\mathrm{mm}$), so the accumulated error leaves history in the solution, and a clean reference requires the strict test from $t = 0$ (solids4foam issue \#552).
A steadiness check on the QoI --- for example a change below $0.04\%$ over the last second --- would have accepted the drifting solution.
The lesson is general to partitioned codes: a convergence measure normalised by a large total quantity can be met without converging the step, and the error it admits is invisible to checks on the solution itself; the residual reduction at acceptance should be reported.

\todo{The Hessenthaler numbers above were obtained on a tree without solids4foam PR \#546; the finding concerns the convergence test and is not expected to depend on it, but confirm with the A/B check. Evidence: solids4foam branch verification/hessenthaler-phaseI, commit 9f33c42a2, \texttt{hessenthalerFsi/verification/README.md}; solids4foam issue \#552.}

\emph{Iterative floors can limit refinement studies.}
In collapsibleChannel the IQN-ILS/Robin--Neumann difference does not decrease with refinement and, at the finest fluid mesh, is comparable with the remaining discretisation error and the reference uncertainty; in blobInTreacle the time study cannot be continued below $\Delta t = 0.03125\,\mathrm{s}$ because the changes reach the tolerance level.
In both, the iterative error rather than the discretisation bounds what can be verified.

\emph{Internal agreement is not external verification.}
IQN-ILS and Robin--Neumann agree to $0.03\%$ (ringAddedMass frequency), $2.2\times10^{-4}$ (womersleyTube), $0.01\%$ (FSI1), and $0.07\%$ (3dTube QoIs).
This bounds the iterative error relative to these differences and checks the implementations against each other; it says nothing about the shared discretisation error, nor about whether the common limit is the continuum solution.
FSI3 shows that the two need not agree to the coupling tolerance when the Robin--Neumann scheme modifies the discrete interface condition.

\emph{Coupling algorithms can affect solution quality as well as cost.}
In ringAddedMass, IQN-ILS produces a mesh-dependent amplitude loss that Robin--Neumann does not (mechanism unknown); and the cheaper algorithm is case dependent (IQN-ILS in ringAddedMass and FSI3, Robin--Neumann in womersleyTube and 3dTube).

\todo{ESSENTIAL: one or two representative coupling-tolerance sweeps (at least two decades, strict test) at fixed discretisation, e.g.\ ringAddedMass (strong level) and collapsibleChannel, so that the claim ``iterative error $\ll$ discretisation error'' rests on more than the Hron--Turek cases; and the $r_1$-at-acceptance audit of Section~\ref{sec:convergence_method}.}

%%%%%%%%%%%%%%%%%%%%%%%%%%%%%%%%%%%%%%%%%%%%%%%%%%%%%%%%%%%%%%%%%%
\subsection{Defects exposed by verification}\label{sec:defects}
%%%%%%%%%%%%%%%%%%%%%%%%%%%%%%%%%%%%%%%%%%%%%%%%%%%%%%%%%%%%%%%%%%

Four defects or hazards in the code under test were found by the studies of Section~\ref{sec:benchmark_suite}, none by comparison with a published value (Table~\ref{tab:defects}).

\begin{table}[htbp]
	\centering
	\footnotesize
	\setlength{\tabcolsep}{3pt}
	\caption{Defects and hazards exposed by verification, and how they were found.}
	\label{tab:defects}
	\begin{tabular}{p{2.6cm}p{5.6cm}p{5.6cm}}
		\toprule
		Case & Defect & How it was found \\
		\midrule
		womersleyTube & $O(\Delta t)$ inconsistency of the end flux under a mixed velocity condition & temporal order falling towards one against an exact reference \\
		3dTube & wall viscous force miscompiled at \texttt{-O3} (temporary-storage aliasing under \texttt{\_\_restrict\_\_}) & deterministic $1.5$--$3\%$ difference in one QoI between two platforms \\
		3dTube & degenerate barycentric weights below a fixed absolute area threshold; interpolation operator changed with mesh level & first divergence after the fix above; broken symmetry and rank dependence \\
		Hessenthaler & legacy coupling test $\min(r_1, r_2)$ accepts unconverged steps in quasi-steady phases & slow drift larger than the spatial differences; residuals at acceptance \\
		\bottomrule
	\end{tabular}
\end{table}

Two observations follow.
First, the toolchain belongs in the verification matrix.
The 3dTube aliasing defect produced a reproducible difference between two platforms that looked like a platform or round-off effect, and a study on one platform, at any resolution, would have reported either the correct or the incorrect result without warning.
Its root cause lies in the reuse of temporary storage by OpenFOAM's field products combined with \texttt{\_\_restrict\_\_}, so it is not specific to solids4foam; the audit for the fix found about 15 further sites of the same mechanism in solids4foam.
Second, a defect can switch on with refinement: the barycentric degeneracy affected levels 2 and 3 of 3dTube but not level 1, so the mesh study compared different interpolation operators, and its apparent convergence of the peak displacement was not evidence about the discretisation.
Cross-platform replication, decomposition invariance and symmetry checks are inexpensive and were decisive; they should accompany refinement studies of coupled codes.

%%%%%%%%%%%%%%%%%%%%%%%%%%%%%%%%%%%%%%%%%%%%%%%%%%%%%%%%%%%%%%%%%%
\subsection{Reliability of published benchmark values}\label{sec:published_reliability}
%%%%%%%%%%%%%%%%%%%%%%%%%%%%%%%%%%%%%%%%%%%%%%%%%%%%%%%%%%%%%%%%%%

Among the literature-referenced problems in the suite, only the Featflow FSI1 table and the regenerated collapsible-channel reference are both independent and accompanied by convergence evidence at the resolution used here (Table~\ref{tab:published}).
This is a statement about the specific values used, not about the quality of the FSI benchmark literature in general.

\begin{table}[htbp]
	\centering
	\footnotesize
	\setlength{\tabcolsep}{3pt}
	\caption{Issues found with published benchmark values.}
	\label{tab:published}
	\begin{tabular}{p{2.8cm}p{5.8cm}p{5.0cm}}
		\toprule
		Case & Issue & Consequence \\
		\midrule
		Hron--Turek FSI2/3 & Summary values of~\citet{Turek2006} differ from Featflow tables (FSI3 drag amplitude $-18\%$; rounded frequencies); Featflow level 3$\to$4 changes 1.6--4.5\% (FSI3 displacements and force amplitudes), 0.9--2.5\% (FSI2) & Choice of source changes conclusions; FSI3 reference uncertain by 3--5\% in $u_x$ and force amplitudes \\
		3dTube & Independent references use implicit Euler, $\Delta t = 10^{-4}\,\mathrm{s}$, without time-step studies; peaks 20--25\% below BDF2 values; figure-digitised & Implicit Euler in the present code reproduces 92--115\% of the gap \\
		mokLidDrivenCavity & Two inconsistent problem definitions; mesh-dependent openings; no convincing convergence studies; group-B spread 5.7\% (peak) & No published value can be treated as correct \\
		collapsibleChannel & Published problem (pre-stressed, massless wall) unsuitable for continuum/partitioned comparison & Modified problem and regenerated reference used \\
		beamInCrossFlow & Quoted values ($5.95\times10^{-5}\,\mathrm{m}$, $1.33\,\mathrm{N}$) are rounded extrapolations of \citet{Richter2012}, often cited to a later paper; earlier calculations solved a different beam position, monitoring point and inflow & Comparison withdrawn; campaign rebuilt to the published definition \\
		cavityFlexibleBottom & Same-lineage, digitised & Offset of $\sim5\%$ unexplained \\
		blobInTreacle & Preprint figures only; first-order time, coarse steady mesh & Steady offset of 7--20\% unexplained \\
		\bottomrule
	\end{tabular}
\end{table}

The table illustrates the central distinction of Section~\ref{sec:verification_methodology}: a widely used benchmark problem does not automatically provide a trustworthy benchmark solution.
The same problem can carry several published ``reference values'', and an $18\%$ difference in a quoted force amplitude exceeds the 2x-to-4x change of every FSI3 quantity.
Disagreements between codes can originate in the discretisation of the reference rather than of the code under test (3dTube), so a discrepancy with a published curve should not be resolved by tuning the code; and same-lineage references, though useful for regression and for diagnosing set-up differences, cannot expose common-mode errors.

%%%%%%%%%%%%%%%%%%%%%%%%%%%%%%%%%%%%%%%%%%%%%%%%%%%%%%%%%%%%%%%%%%
\subsection{Cost and dimensionality}\label{sec:cost}
%%%%%%%%%%%%%%%%%%%%%%%%%%%%%%%%%%%%%%%%%%%%%%%%%%%%%%%%%%%%%%%%%%

\begin{table}[htbp]
	\centering
	\footnotesize
	\setlength{\tabcolsep}{3pt}
	\caption{Recorded computational cost of the verification studies. Hardware and core counts differ between cases, so the figures indicate orders of magnitude only.}
	\label{tab:cost}
	\begin{tabular}{p{3.0cm}p{10.8cm}}
		\toprule
		Case & Recorded cost \\
		\midrule
		ringAddedMass & tutorial $\approx75$\,s serial; full study (35 runs) $\approx30$\,min with 32 concurrent serial runs \\
		womersleyTube & runs 2.5--25\,min serial; full study $\approx30$\,min on 5 cores \\
		collapsibleChannel & tutorial $\approx10$\,min; finest fluid mesh 70\,min; full sweep $\approx2.5$\,h (one case per core); serial only \\
		Hron--Turek FSI1 & 1x 781\,s serial; 2x 3\,062\,s on 2 ranks; 4x estimated $>14$\,h on 3 ranks \\
		Hron--Turek FSI3 & 1x 1.4\,h serial; 2x 5.5\,h on 8 ranks; 4x 40.8\,h on 32 ranks (no speed-up beyond 32 ranks); OpenFOAM-v2512 \\
		Hron--Turek FSI2 & 1x 1.1--1.2\,h serial; 2x 4.0--4.1\,h on 7--8 ranks \\
		3dTube & level 1 504\,s on 4 ranks; level 2 3\,641\,s on 8 ranks; level 3 22\,502\,s (6.25\,h) on 32 ranks ($\approx720$ core-hours) \\
		mokLidDrivenCavity & $96^2$ 3\,729\,s on 2 ranks; quick run $\approx7$\,min \\
		cavityFlexibleBottom$^\ast$ & three-level sweep $\approx10$\,min serial \\
		beamInCrossFlow$^\ast$ & in progress \\
		blobInTreacle$^\ast$ & full study $\approx33$\,min \\
		\bottomrule
	\end{tabular}
\end{table}

Cost rises by roughly four orders of magnitude from the exact-reference problems to the finest three-dimensional calculations (Table~\ref{tab:cost}), and the completeness of the evidence falls in the opposite direction: the inexpensive two-dimensional and axisymmetric problems have three or more levels in space and time and exact or independent references, whereas the three-dimensional problems have three levels with sub-nominal or undefined orders (3dTube) or, so far, no completed study (beamInCrossFlow, Hessenthaler). Even in two dimensions a third level can cost an order of magnitude more than the second (Hron--Turek FSI3: 40.8\,h on 32 ranks at 4x, without further strong scaling).
This illustrates one practical reason why rigorous convergence evidence becomes scarcer as geometric and computational complexity increase, and it motivates a tiered suite (Section~\ref{sec:recommended_suite}).
Which error source dominates is case dependent: in the thin-walled two-dimensional collapsibleChannel the solid discretisation can dominate ($17.4\%$ versus $0.44\%$ for linear versus cubic solid), whereas in Hron--Turek FSI3 refining the solid alone accounts for only $\approx5\%$ of the change in the $u_y$ amplitude.
\todo{Revisit once the beamInCrossFlow and Hessenthaler campaigns are complete; add cost per degree of freedom and per time step if comparable hardware is used.}
